% SOURCE_AUTHORITY: sga5_fr_workpass.tex, lines 13017--13048 (LF-normalized unit).
% SOURCE_SHA256: 791F4EFFC5E02832D5D77ED03518C8156D6F07E4C8238B03545DB93D883FBB28
Por outro lado, tem-se o triângulo distinguido (onde $Y=C-U$)
\[
\begin{tikzcd}[column sep=huge,row sep=large]
& \RG_Y(F|Y) \arrow[dl] & \\
\RG_C(F_{U,C}) \arrow[rr] && \RG_C(F) \arrow[ul]
\end{tikzcd}
\tag{7.8}
\]
onde, evidentemente,
\[
\RG_Y(F)=\coprod_{y\in Y}F_y\in\Ob(\Delta),
\]
uma vez que $F_y\in\Ob(\Delta)$ (condição (i)). Resulta, portanto, que $\RG_C(F)\in\Ob(\Delta)$, ou seja, a primeira afirmação do teorema. O triângulo distinguido (7.8) fornece, além disso, a relação
\[
\clK{\Delta}(\RG_C(F))
=
\clK{\Delta}(\RG_C(F_{U,C}))
+\sum_{y\in Y}\clK{\Delta}(F_y),
\]
e, utilizando o isomorfismo 7.7, obtém-se:
\[
\clK{\Delta}(\RG_C(F))
=
\clK{\Delta}\bigl(\RHom_{\Lambda[G]}(\RG_{C'}(\Lambda_{U',C'})^\vee,F_{\bar\eta})\bigr)
+\sum_{y\in Y}\clK{\Delta}(F_y).
\]
Para obter a fórmula de Euler--Poincaré basta agora utilizar a fórmula de Weil (5.2). Com efeito, primeiro 4.4, (5.2) e 3.6 (no caso particular $B=\Lambda[G]$) fornecem
\[
\clKstar{\Lambda[G]}\bigl(\RG_{C'}(\Lambda_{U',C'})\bigr)
=
\clKstar{\Lambda[G]}\bigl(\RG_{C'}(\Lambda_{U',C'})^\vee\bigr).
\]
